% Additive insertion only. Do not delete the reviewed-simulation premise,
% Horn/event-normalization caveats, or conditional transfer statements.
\subsection{A Restricted Ground Scheduler Lemma}
\label{app:restricted-scheduler}

The following result concerns two explicitly defined scheduler reductions,
not the complete symbolic semantics of ProVerif and Tamarin. It does not
discharge the reviewed process-to-rule simulation assumed in the proposed
inclusion chain.

Fix exactly one completed initialization and its identity valuation.
The generated private state channel and initial fresh identity remain;
there is no fresh allocation during updates. Admit only the linear
ProVerif worker and the standard Tamarin rule shape, an acyclic finite
state graph, and terms ground after this valuation. Require a free term
algebra with pair projections, hygienic identifiers, parseable generated
texts, pairwise disjoint dispatch patterns at each state, and at most one
observed event per case. Exclude tables, responders, leaks, persistent
state, Clone, and Rollback. Initial knowledge is the deductive closure of
public constants and the actual initialization output; in particular,
the fixture's identity is public by tuple projection. These conditions
are additional restrictions, not certificates issued by the current checker.

Write a case as $(s,p,e,[o_1,\ldots,o_r],t)$, where $p$ is a ground
input or absent and $e$ is one observed event or the empty word.
Let $D$ denote public deductive closure. Define the atomic scheduler by
\[
 (s,K,w)\longrightarrow
 (t,D(K\cup\{o_1,\ldots,o_r\}),we),
\]
when $p\in D(K)$, or when $p$ is absent.
Define the worker scheduler to reserve the unique private state output,
receive a public dispatch, reduce all parallel pattern and state tests,
and either abandon the state or select the unique successful case.
The successful case emits $e$, delivers the public outputs in order,
and finally releases the private successor. With no observed event,
commit is a silent point after successful checks and before outputs.
An absent case input still requires a public dispatch in the printed
worker; the public numeral constant supplies one.

\noindent\textbf{Restricted scheduler lemma.}
These two scheduler reductions have the same finite observed event words.

\emph{Proof.}
For the forward direction relate available, reserved, dispatched,
selected, and abandoned worker states at $s$ to the atomic state
$(s,D(K),w)$. After commit of case $c$, if exactly $j$ outputs have
been delivered, relate the worker instead to
\[
 (t,D(K\cup\{o_{j+1},\ldots,o_r\}),w).
\]
Reservation, dispatch, failed pattern tests, failed state tests, and
inactive sibling reductions are silent matches. Disjointness ensures
that at most one parallel branch succeeds. If none succeeds, the worker
abandons the token and the atomic scheduler simply retains it without
being scheduled again. Before commit there is no output debt, so the
successful input is justified by current, not future, knowledge.
Commit is matched by the same atomic case, with either one observation
or none. Each public delivery is a silent match because
\[
 D(D(K\cup\{o_{j+1}\})\cup\{o_{j+2},\ldots,o_r\})
 =D(K\cup\{o_{j+1},\ldots,o_r\}).
\]
Once the debt is empty, private successor release is silent. Replication
unfolding and extra workers without a token add no observation. An
abandoned or terminal worker adds none. Stopping at any partial prefix,
including immediately after the event, is admissible. Induction over
worker reductions proves forward trace inclusion.

For the reverse direction use a different relation: a stable available
worker and an atomic state must have exactly the same state, knowledge,
and word, without debt. For an enabled atomic case, reserve the token,
send its deducible input (or a public constant for absent input), reduce
failed alternatives, select the unique successful branch, emit its zero
or one observed event, deliver all outputs, and release the successor.
This finite path restores the stable relation. Zero outputs require
only release; multiple outputs are serialized before the next token
can be consumed. Empty runs and terminal prefixes require no step.
Induction over atomic runs proves reverse trace inclusion.
\hfill$\square$

This is not reverse simulation from every forward-related configuration:
after a bad dispatch an abandoned worker cannot match the atomic
scheduler's next event. It is not branching bisimulation or a liveness
result. The public-knowledge sets can differ at the same observed prefix.
Semantic adequacy of these scheduler reductions for the full native
languages, including attacker processes and structural congruence,
remains an explicit obligation.

\paragraph{A generated-code boundary.}
An independent negative control publishes $(I,\mathsf{seal}(n_1))$,
starts at $n_0$, and gives its sole case input
$(\mathsf{a},\mathsf{seal}(x))$. The constructor is private.
The Tamarin rule can bind $x=n_1$ and emit the event. The ProVerif
printer instead emits a constructor equality pattern containing its
already bound internal state variable \texttt{x}, requiring
$\mathsf{seal}(n_0)$, which is not available. The existing checker
reports no FAIL. On the generated files ProVerif proves the event
unreachable while Tamarin verifies its reachability. Closure of
constructor equality terms and capture-avoiding identifiers are therefore
additional necessary guards for a broader claim. A separate scalar-input
negative control fails ProVerif parsing despite a passing report.

The independent suite passes 21 tests. A finite ground fixture has three
atomic states, thirteen worker states (three abandoned), and three event
words on each side. Native reachability checks support the fixture and
the negative controls, not full trace equivalence. The previous two-thread
bounded experiment remains a separate finite abstraction. Neither result
removes the original unproved simulation and Horn-abstraction assumptions,
nor upgrades transfer-dependent Achievement Levels.
